\chapter{Installation et prise en main}

\section{Installer l'application}
\soft est distribué sous forme de fichier APK. Copiez-le sur votre téléphone ou téléchargez-le, puis ouvrez-le.
Android vous demande d'autoriser l'installation depuis cette source : acceptez pour cette fois. Les mises à jour
s'installent de la même façon, par-dessus la version précédente ; vos listes et vos réglages sont conservés.

\section{Les autorisations}
Au premier lancement, \soft vous demande l'accès à vos musiques. Sur les versions récentes d'Android, vous
choisissez un dossier précis plutôt que de donner accès à tout le téléphone ; nous vous conseillons de désigner le
dossier qui contient votre collection. L'autorisation des notifications sert à afficher la lecture en cours et la
progression quand vous ajoutez un dossier entier à une liste.

\section{Les trois écrans}
L'application tient en trois écrans, accessibles par la barre du bas. L'écran \textbf{Lecture} montre le morceau en
cours, l'oscilloscope, le tableau de bord de la puce et les commandes. L'écran \textbf{Listes} rassemble vos listes
de lecture. L'écran \textbf{Navigateur} donne accès à vos fichiers et aux catalogues en ligne.

\capturedouble{now-playing-en}{playlist-en}{Écran Lecture}{écran Listes}

Pour votre première écoute, ouvrez le navigateur, choisissez un catalogue (la High Voltage SID Collection est un bon
point de départ) et touchez un morceau. La lecture démarre et l'écran Lecture se remplit : en haut l'oscilloscope
de chaque voix, en dessous le tableau de bord de la puce. Si vous préférez l'écran épuré, repliez les informations
avec la flèche en bas à gauche.

\capture{browser-sources-en}{Navigateur : les sources}
